\documentclass[aspectratio=169]{beamer}
\mode<presentation>

\usepackage[UTF8,fontset=fandol]{ctex}
\usepackage{pku_beamer}

\PKUSetupCJKFonts

\title[Short Title]{Presentation Title}
\subtitle{Presentation Subtitle}
\author[Author Name]{Author Name}
\institute[PKU SEE]{School of Environment and Energy, Peking University}
\date{\today}

\begin{document}

\PKUTitlePage
\PKUOutlineFrame

\section{Background}

\begin{frame}{Motivation}
  \begin{itemize}
    \item State the research question or presentation context.
    \item Use \PKUHighlight{highlighted keywords} for important terms.
    \item Keep each slide focused on one main message.
  \end{itemize}
\end{frame}

\begin{frame}{Problem Statement}
  \begin{block}{Core question}
    What problem is being solved, and why does it matter?
  \end{block}

  \vspace{0.35cm}

  A concise formulation can be written as
  \[
    \hat{\theta} = \arg\min_{\theta} \mathcal{L}(\theta; X, y).
  \]
\end{frame}

\section{Method}

\begin{frame}{Workflow}
  \begin{enumerate}
    \item Prepare the input data.
    \item Define the model or analytical method.
    \item Evaluate the result with clear metrics.
  \end{enumerate}
\end{frame}

\begin{frame}{Result Summary}
  \begin{table}
    \centering
    \caption{Example table}
    \begin{tabular}{lcc}
      \toprule
      Method & Metric A & Metric B \\
      \midrule
      Baseline & 0.71 & 0.68 \\
      Proposed & 0.84 & 0.79 \\
      \bottomrule
    \end{tabular}
  \end{table}
\end{frame}

\section{Conclusion}

\begin{frame}{Key Takeaways}
  \begin{itemize}
    \item Summarize the most important finding.
    \item Explain the practical implication.
    \item End with one concrete next step.
  \end{itemize}
\end{frame}

\PKUAcknowledgementsFrame

\end{document}
